\scanpage{19}\mspage{45}
\deleted{en groupes \deleted{localement}\inserted{lisse et} de type fini sur $k$, soit $\underline{\mathrm{Pic}}_U$, et un isomorphisme
\[
\mathrm{Pic}(U)=\underline{\mathrm{Pic}}_U(k) \tag*{(8.13)}
\]
de $\mathrm{Pic}(U)$ avec le groupe des points à valeurs dans $k$ de $\underline{\mathrm{Pic}}_U$. Cela nous fournit donc pour $\mathrm{Pic}(U)$ une structure d'extension
\[
0\handwritten{\to}\mathrm{Pic}(U)^{\unsure{o|0}}\handwritten{\longrightarrow}\mathrm{Pic}(U)\handwritten{\longrightarrow}\mathrm{NS}(U)\handwritten{\longrightarrow}0\ , \tag*{(8.14)}
\]
où on a posé
\[
\mathrm{Pic}(U)^o=\underline{\mathrm{Pic}}_U^o(k)\ ,\quad\mathrm{NS}(U)=\underline{\mathrm{NS}}_U(k)\ , \tag*{(8.15)}
\]
\deleted{\unsure{où}} $\underline{\mathrm{Pic}}_U^o$ désignant la composante neutre de $\underline{\mathrm{Pic}}_U$, et $\underline{\mathrm{NS}}_U$ le quotient de $\underline{\mathrm{\deleted{\unsure{B}}Pic}}_U$ par cette composante. La construction indiquée dans loc. cit. permet d'ailleurs de préciser la structure de $\mathrm{NS}(U)$, en introduisant la matrice d'intersection
\[
(\mathrm C_i.\mathrm C_j)_{1\leq i,j\leq r}
\]
définie par les $r$ composantes irréductibles $\mathrm C_i$ de $f^{-1}(x)_{\text{réd}}$ sur le \margin{\handwritten{$\times$}}schéma régulier $X'$ : en vertu de MUMFORD\inserted{\handwritten{[28]}}, cette matrice est définie néggtive, donc le conoyau de l'homomorphisme
\[
\underline{\mathrm Z}^{r}\handwritten{\longrightarrow}\underline{\mathrm Z}^{r}
\]
qu'elle définit est un groupe fini, dont l'o\deleted{\unsure{dr}}rdre n'est autre que la valeur abso\deleted{\unsure{h}}lue du déterminant de cette matrice. C\deleted{\unsure{o}}eci posé, on trouve que $\mathrm{NS}(U)$ est canoniquement isomorphe à ce groupe fini. La matrice envisagée étant symétrique, on voit aussitôt que ce groupe est canoniquement autodual (par un accouplement \deleted{\unsure{parfait}}\handwritten{symétrique} à valeurs dans $\underline{\mathrm Q}/\underline{\mathrm Z}$), d'où une autodualité\inserted{\handwritten{ symétrique}}
\[
\mathrm{NS}(U)\mathbin{\mathrm x}\mathrm{NS}(U)\handwritten{\longrightarrow}\underline{\mathrm Q}/\underline{\mathrm Z}\ , \tag*{(8.16)}
\]
dont on constate aisément qu\deleted{e}'elle est indépendante de la résolution $X'$ choisie.

Si maintenant \handwritten{$\unsure{l|\ell}$} est un nombre premier distinct de $\mathrm{car}\ k$, alors la structure connue des groupes algébriques commutatifs connexes sur $k$ \margin{\handwritten{$\times$}}\inserted{\handwritten{[9, exp.4]}} mont\deleted{\unsure{e}}re que $\mathrm{Pic}(U)\mathrm{o}$ est un groupe \handwritten{$l$}-divisible, donc}
% archive sheet number: 18
